\chapter{Appendix B: Neuron Model Equations}\label{app:neuron_equations}

This appendix documents the neuron and synapse models used by the synthetic nervous systems in this project, both in the AnimatLab walker model and in the SNS-Toolbox implementation of the MuJoCo pipeline (Chapter~\ref{ch:futurework}). The authoritative model state is the working AnimatLab project file (\texttt{Biped\_2xCPG\_wSubs/Biped\_2xCPG\_wSubs.aproj}); after the network reorganization described in Section~\ref{sec:preliminary_sim_methods}, a fresh standalone export (\texttt{Biped\_2xCPG\_wSubs\_Standalone.asim}) runs headless and carries the rhythm-generator charts, and the earlier phase-1 instrumented export (\texttt{walk new new tester added 2 axis\_phase1.asim}) records the run reported in Section~\ref{sec:preliminary_sim_results}.

\section{Spiking Neuron Model}

Each spiking neuron is a conductance-based leaky integrator with threshold and adaptation,
\begin{equation}
    C_{m}\,\dot{V} \;=\; -\,g_{\text{leak}}\left(V - E_{\text{rest}}\right)
                      \;+\; I_{\text{syn}} \;+\; I_{\text{app}},
    \label{eq:app_membrane}
\end{equation}
where $V$ is the membrane potential, $C_{m}$ the membrane capacitance ($g_{\text{leak}} = C_{m}/\tau_{m}$ with membrane time constant $\tau_{m}$), $E_{\text{rest}}$ the resting potential, $I_{\text{syn}}$ the summed synaptic current, and $I_{\text{app}}$ any applied (tonic stimulation) current. A spike is emitted when $V$ exceeds the firing threshold, after which $V$ resets and an after-hyperpolarization (AHP) of amplitude $A_{\text{AHP}}$ with time constant $\tau_{\text{AHP}}$ is applied; the threshold itself accommodates with the firing history (relative accommodation $r_{\text{acc}}$, time constant $\tau_{\text{acc}}$). Optional calcium channels are defined but disabled in the current model ($G_{\text{max,Ca}} = 0$).

The parameter values shared by the neurons in the walker model are given in Table~\ref{tab:app_neuron}.

\begin{table}[htbp]
\caption{Spiking neuron parameters in the AnimatLab bipedal walker model
(\texttt{Biped\_2xCPG\_wSubs}). All neurons share these values except the motor
neuron threshold noted below.}
\label{tab:app_neuron}
\begin{tabular}{@{}lll@{}}
\toprule
\textbf{Parameter} & \textbf{Symbol} & \textbf{Value} \\
\midrule
Resting potential        & $E_{\text{rest}}$   & \qty{-60}{\mV} \\
Membrane time constant   & $\tau_{m}$          & \qty{5}{\ms}   \\
Initial firing threshold & $V_{\text{th}}$     & \qty{-55}{\mV} (motor neurons: \qty{-59}{\mV}) \\
Relative accommodation   & $r_{\text{acc}}$    & 0.3            \\
Accommodation time const & $\tau_{\text{acc}}$ & \qty{10}{\ms}  \\
AHP amplitude            & $A_{\text{AHP}}$    & \qty{1}{\mV}   \\
AHP time constant        & $\tau_{\text{AHP}}$ & \qty{3}{\ms}   \\
\bottomrule
\end{tabular}
\end{table}

\section{Non-Spiking (Graded) Chemical Synapses}

Pattern-formation and motoneuron-drive connections in the walker use graded chemical synapses, whose conductance is a clipped linear ramp of the presynaptic potential,
\begin{equation}
    g_{\text{syn}}(V_{\text{pre}})
    \;=\; g_{\text{max}}\;
    \operatorname{clamp}\!\left(
        \frac{V_{\text{pre}} - V_{\text{th}}}{V_{\text{sat}} - V_{\text{th}}},
        \,0,\,1\right),
    \qquad
    I_{\text{syn}} = g_{\text{syn}}\left(V_{\text{post}} - E_{\text{syn}}\right),
    \label{eq:app_nonsyn}
\end{equation}
where $V_{\text{th}}$ and $V_{\text{sat}}$ are the synaptic threshold and saturation potentials and $E_{\text{syn}}$ the reversal potential. The walker model instantiates three non-spiking synapse types:

\begin{itemize}
    \item \textbf{Nicotinic ACh (excitatory drive):}
          $E_{\text{syn}} = \qty{-10}{\mV}$,
          $g_{\text{max}} = \qty{0.5}{\micro\siemens}$,
          $V_{\text{th}} \to V_{\text{sat}} = \qty{-65}{\mV} \to \qty{-20}{\mV}$.
    \item \textbf{Hyperpolarizing IPSP (inhibition):}
          $E_{\text{syn}} = \qty{-70}{\mV}$,
          $V_{\text{th}} \to V_{\text{sat}} = \qty{-55}{\mV} \to \qty{-40}{\mV}$.
    \item \textbf{Depolarizing IPSP:}
          $E_{\text{syn}} = \qty{-55}{\mV}$,
          $V_{\text{th}} \to V_{\text{sat}} = \qty{-60}{\mV} \to \qty{-40}{\mV}$.
\end{itemize}

Spiking chemical synapses and electrical synapses (rectifying coupling weight 0.1, non-rectifying 0.2) are also defined in the model and are available for the sensory afferent pathways of the MuJoCo/SNS implementation.

\section{Relation to the SNS-Toolbox Implementation}

These are the standard AnimatLab firing-rate/conductance models. SNS-Toolbox \citep{nourse_sns_2023} implements corresponding classes of non-spiking and spiking neurons with conductance-based synapses, and the AnimatLab parameters above served as the physically motivated initialization for the first SNS-Toolbox instantiation of the spinal network: a 406-neuron network with per-leg rhythm-generator half-centers, phase-shifted pattern-formation groups, motoneuron pools for all 92 muscle-tendon actuators of the converted model, and Ia, Ib, and II afferent populations. In that implementation the non-spiking membrane and graded-synapse equations take the same forms as Equations~(\ref{eq:app_membrane}) and~(\ref{eq:app_nonsyn}) with the spiking threshold and accommodation terms removed, and the afferent resting-tone offsets are zeroed so that Ia and Ib signals are pure activity measures; the rhythm-generator layer has been verified to oscillate at a human-like period with the two legs in antiphase under tonic drive (Section~\ref{sec:fw_cpg}). An explicit parameter-mapping and validation campaign between the two implementations remains future work; the tonic stimulation tests of Section~\ref{sec:fw_cpg} correspond to constant $I_{\text{app}}$ injections into Equation~(\ref{eq:app_membrane}) at the motoneuron, pattern-formation, and rhythm-generator layers respectively.

